\documentclass[11pt,a4paper]{article}
\usepackage[T1]{fontenc}
\usepackage[utf8]{inputenc}
\usepackage{amsmath,amssymb,xcolor}
\usepackage[margin=25mm]{geometry}
\usepackage[colorlinks=true,urlcolor=blue]{hyperref}
\definecolor{nexorablue}{HTML}{2459B5}
\setlength{\parindent}{0pt}
\setlength{\parskip}{7pt}
\setlength{\emergencystretch}{2em}
\renewcommand{\familydefault}{\sfdefault}
\begin{document}

{\small\bfseries\color{nexorablue} NEXORA CPGE}\par
{\small MP / KUT-SPE-07-01 / Énoncé}\par
{\LARGE\bfseries Diagonalisation avec une valeur propre double}\par
\vspace{0.25em}\hrule\vspace{0.6em}

Pour la matrice ci-dessous, trouver une matrice inversible $P$ et une matrice diagonale $D$ telles que $P^{-1}AP=D$.
\[ A=\begin{pmatrix}2&0&0\\1&4&-1\\-2&-4&4\end{pmatrix}. \]

\bigskip
{\footnotesize Adaptation française par Nexora CPGE. Source : Ken Kuttler (Lyryx),
\href{https://math.libretexts.org/Bookshelves/Linear_Algebra/A_First_Course_in_Linear_Algebra_(Kuttler)/07\%253A_Spectral_Theory/7.02\%253A_Diagonalization}{\textit{A First Course in Linear Algebra}}, § 7.2, exemple 1.
\href{https://creativecommons.org/licenses/by/4.0/}{CC BY 4.0}.}
\end{document}
